% Generated draft: requires mathematical and source review.
\documentclass[11pt]{article}
\usepackage[margin=25mm]{geometry}
\usepackage{fontspec}
\setmainfont{Latin Modern Roman}
\usepackage{amsmath,amssymb,mathtools,unicode-math}
\setmathfont{Latin Modern Math}
\usepackage{enumitem,needspace,xurl,hyperref}
\setlength{\parindent}{0pt}
\setlength{\parskip}{5pt}
\setlength{\emergencystretch}{3em}
\newcommand{\N}{\mathbb{N}}
\newcommand{\Z}{\mathbb{Z}}
\newcommand{\Q}{\mathbb{Q}}
\newcommand{\R}{\mathbb{R}}
\newcommand{\C}{\mathbb{C}}
\newcommand{\F}{\mathbb{F}}
\newcommand{\A}{\mathbb{A}}
\newcommand{\PP}{\mathbb{P}}
\newcommand{\E}{\mathbb{E}}
\newcommand{\eps}{\varepsilon}
\newcommand{\dd}{\,\mathrm{d}}
\begin{document}
\section*{UnsolvedMath ID: 30005421}
\textbf{Source problem code:} catalogue entry
\subsection{Definitions and mathematical statement}
A quiver is a directed graph \(Q=(Q_0,Q_1,s,t)\), where \(Q_0\) is its vertex set, \(Q_1\) its arrow set, and \(s,t:Q_1\to Q_0\) give the source and target of each arrow. The supplied record asks the following conjectural question:
\[
\text{For every wild quiver }Q,\quad
\exists\,\mu\text{ a stability condition on }Q
\quad\text{such that }(Q,\mu)\text{ has uniform expansion.}
\]
Here ``wild'', ``stability condition'', and ``uniform expansion'' retain their intended meanings in the cited report. Their definitions are absent from the supplied excerpt. In particular, the record does not specify the admissible quivers, the base field for their representations, the data defining \(\mu\), or the constants, representation families, subobjects, and quantifier order required by uniform expansion. Consequently, this is a faithful statement of the supplied target, but a complete intrinsic formulation cannot be recovered from the supplied evidence. No finiteness, connectedness, or acyclicity hypothesis on \(Q\) is imposed by this draft, and their possible necessity remains unresolved.

The accompanying dimension-expander discussion explicitly takes \(k\) to be algebraically closed and \(V\) to be a finite-dimensional \(k\)-vector space. It discusses operators \(F,G\in\operatorname{End}_k(V)\) and dimensions of \(U+F(U)+G(U)\). It describes the existence of simple representations of the free algebra in two generators in arbitrary finite dimensions through the condition
\[
\dim_k\bigl(U+F(U)+G(U)\bigr)>\dim_k U
\qquad\text{for every }0\ne U\subsetneq V.
\]
It also states a stronger quantitative expansion assertion for nonzero subspaces subject to a dimension bound. That bound and the numerical expansion factor are corrupted in the supplied extraction and are not reconstructed here. The excerpt attributes the qualitative quantitative assertion to its reference [1], and a determination of an asserted optimal constant to its reference [2]. It describes the quantitative property as Zariski-open in the entries of \(F,G\), hence holding for almost all pairs, and the proof in [2] as completely nonconstructive. These are supplied contextual assertions, not independently verified results in this draft. This example does not supply the missing definition of uniform expansion for \((Q,\mu)\), nor establish that its field hypothesis governs the conjecture.
\subsection{Short English statement}
The supplied conjecture asks whether every wild quiver admits a stability condition giving uniform expansion. The excerpt omits the definitions and hypotheses needed to make this fully precise. Its open status and verification labels are supplied curation claims; this statement-only draft requires human mathematical review.
\subsection{Sources}
\begin{itemize}
\item Local UnsolvedMath registry: supplied record with permanent ID \(30005421\), title \emph{Uniform Expansion Stability for Wild Quivers}, exact source problem code \texttt{catalogue entry}, in the Oberwolfach Reports Open Problems collection. The record gives proposed year 2023 and identifies no individual proposer.
\item \emph{Representation Theory of Quivers and Finite-Dimensional Algebras} (2023), Oberwolfach Reports, pp.\ 397--486, DOI 10.4171/owr/2023/7. Supplied source report and literature-source link: \url{https://doi.org/10.4171/owr/2023/7}. The excerpt says ``We conjecture'' without identifying the speaker in the supplied record.
\item Supplied later literature reference: \emph{J. Comb. Algebra} \textbf{8} (2024), 111--119, \url{https://ems.press/content/serial-article-files/47504}. No article title or authors are supplied.
\item Literature assessment checked at \texttt{2026-08-22}: supplied dated curation, not independently verified here. It labels the problem open, says that the source report states the wild-quiver converse as a conjecture, and reports that the later paper proves broad special classes while the searched literature does not establish the assertion for every wild quiver. The special classes and search scope are not specified. The record expressly requests independent verification of time-sensitive status and cited claims.
\item The supplied statement status is \texttt{corrected\_verified}; its verification note claims comparison with the official report and preservation of hypotheses, notation, and target. This draft records that claim without independently confirming it. The excerpt's references [1] and [2] have no supplied bibliographic entries and cannot be identified from this record.
\end{itemize}
\end{document}
